% ---
% PREÂMBULO LaTeX - Organizado e Conciso
% ---

% --- Codificação, Fonte e Idioma ---
\usepackage{inputenc}           % Codificação
% \usepackage{fontenc}            % Codificação de fontes
\usepackage{lmodern}            % Fonte Latin Modern
\usepackage{microtype}          % Melhorias tipográficas (justificação, espaçamento)

% --- Layout e Estrutura ---
\usepackage{lastpage}           % Referência à última página (ficha catalográfica)
\usepackage{indentfirst}        % Indenta primeiro parágrafo das seções

% --- Elementos Visuais (Cores, Gráficos) ---
\usepackage{xcolor}             % Controle avançado de cores
\usepackage{graphicx}           % Inclusão de gráficos

% --- Matemática ---
\usepackage{amsmath}            % Recursos matemáticos avançados
\usepackage{amssymb}            % Símbolos matemáticos adicionais
\usepackage{amsfonts}           % Conjuntos de fontes matemáticas

% --- Listagens de Código ---
\usepackage{listings}           % Blocos de código fonte

% --- Citações e Referências ---
% \usepackage[brazilian,hyperpageref]{backref} % Referências de página na bibliografia (opcional)

% --- Ferramentas de Edição e Anotação ---
\usepackage{cancel}             % Riscar texto (\cancel{})
\usepackage{soulutf8}           % Sublinhar, destacar, etc. (\st{}, \hl{})
\setstcolor{red}                % Cor para \st{} (strikeout)

% ---
% CONFIGURAÇÕES GLOBAIS E DE PACOTES
% ---
\definecolor{blue}{RGB}{41,5,195} % Azul personalizado
\colorlet{oliveGreen}{blue!20!black!50!green} % Verde oliva para anotações

\makeatletter % Permite uso de @ em comandos (para hyperref e outros)

% Configurações do Hyperref (links e metadados do PDF)
\hypersetup{
    colorlinks=true,
    linkcolor=blue,
    citecolor=blue,
    filecolor=magenta,
    urlcolor=blue,
    bookmarksdepth=4,
    % pdftitle={\@title}, pdfauthor={\@author}, % Metadados (se definidos pela classe)
    % pagebackref=true, % Backreferences (requer pacote backref)
}

% Aparência: Título do Sumário, Indentação do Sumário, Legendas
\renewcommand{\printtoctitle}[1]{\normalsize\bfseries\centering TABLE OF CONTENTS}
% \cftsetindents{chapter}{0em}{0em} % Ajustar conforme estrutura (part/chapter)
\cftsetindents{section}{1.5em}{2.5em}
\cftsetindents{subsection}{3em}{3.5em}
\cftsetindents{subsubsection}{4.5em}{5em}
\AtBeginDocument{\captionsetup{labelfont=bf}} % Legendas em negrito

% --- Comandos Personalizados ---
\newcommand{\cito}[1]{\textsuperscript{\cite{#1}}} % Citação sobrescrita
\newcommand{\PROFESSOR}[1]{\textcolor{red}{\MakeUppercase{#1}}}
\newcommand{\PROFESSORDEL}[1]{\textcolor{red}{\st{#1}}}
\newcommand{\PROFESSORADD}[1]{\textcolor{oliveGreen}{#1}}
\newcommand{\PROFESSORREP}[2]{\PROFESSORDEL{#1}\PROFESSORADD{#2}}

\makeatother % Restaura significado de @

% --- Espaçamentos ---
\setlength{\parindent}{1.3cm}    % Indentação do parágrafo
\setlength{\parskip}{0.2cm}      % Espaçamento entre parágrafos

% --- Índice Remissivo ---
\makeindex % Habilita compilação do índice remissivo
% ---